\section{附录：一份可直接复用的 Benchmark 报告模板}

\subsection{报告模板字段}
为了把课程思想直接落到团队协作，这里给出一份可复用模板。核心原则是 ``所有可影响结论的变量都要显式记录''，避免后续复现时出现信息缺失。

\begin{longtable}{p{2.9cm}p{3.7cm}p{4.5cm}}
\toprule
\textbf{字段} & \textbf{示例} & \textbf{说明} \\
\midrule
Model Version & model\_2026\_04\_rc2 & 必须能唯一定位权重与配置 \\
Benchmark Version & bench\_v3.4 & 任务集与样本切分版本号 \\
Prompt Protocol & proto\_2026\_04\_a & system prompt + template + few-shot 策略 \\
Decode Config & temp=0.2, top\_p=0.95 & 生成参数是评测协议组成部分 \\
Judge Config & judge-v2 & 若有 LLM judge，需记录模型和 rubric \\
Runs / Seeds & 10 runs, seeds=[1..10] & 无重复运行的结果不应进入最终报告 \\
Main Metric & success\_rate=73.0 & 总体分数 \\
Uncertainty & std=1.2, 95\%CI=[72.2,73.8] & 至少报告一种置信区间 \\
Hard Split & 58.9 (+0.1) & 避免总体分数掩盖难样本退化 \\
Reliability Metrics & invalid\_action, timeout, recovery & Agent 任务必须补充可靠性指标 \\
Environment Version & env-2026-04-01 & 保证任务环境可复现 \\
Known Limitations & judge bias risk, sample drift & 主动披露风险与限制 \\
\bottomrule
\end{longtable}

\subsection{发布前检查清单}
\begin{importantbox}{Release Checklist}
\begin{enumerate}
    \item 是否完成了不少于 5 次重复评测并保存全部原始日志？
    \item 是否确认协议版本、judge 版本和环境版本均已冻结？
    \item 是否报告了 hard split 与关键业务子集表现？
    \item 是否有人评抽检来校准自动评审偏差？
    \item 是否有回归告警阈值与回滚预案？
\end{enumerate}
\end{importantbox}

\begin{knowledgebox}{为什么这份模板有用}
它把课程中的统计原则、协议治理和工程可复现性收敛成一份统一文档格式。模板统一后，跨团队比较将基于同一信息面展开，能显著降低 ``因为记录不全导致的争议''。
\end{knowledgebox}

\subsection{本章小结}
评测体系要可持续，必须让 ``报告结构'' 标准化。模板和清单看似琐碎，但它们决定了 benchmark 结论能否被长期复用。

